\section{Two standard inputs and an analytic consequence}
\label{sec:preliminaries}

We first state the two established results used in the reduction.
The new combinatorial and analytic arguments will be proved in full.
For a positive integer $r$, write $[r]=\{1,\ldots,r\}$.
All products of metric spaces carry the sup metric unless stated otherwise.
The circle $\T=\R/\mathbb Z$ carries the distance
$d_{\T}(x,y)=\min_{k\in\mathbb Z}|x-y-k|$.
Uniform inputs on a cube have product Lebesgue law, and uniform inputs on
a product of circles have product Haar law. Norms $\|\cdot\|_p$
always use the indicated probability measure.
An empty sum is zero.

\subsection{A projection constraint problem}

A \emph{Label Cover instance} consists of finite question sets $U,V$,
nonempty finite answer alphabets $\Sigma_L,\Sigma_R$, and a nonempty
explicit list $\mathcal C$ of test occurrences.
An occurrence $c$ has questions $u_c\in U$, $v_c\in V$, and a total map
$\pi_c:\Sigma_L\to\Sigma_R$.
The value is
\[
 \val(\mathcal C)=
 \max_{\ell:U\to\Sigma_L,\ \rho:V\to\Sigma_R}
 \Pr_{c\in\mathcal C}\bigl[\pi_c(\ell(u_c))=\rho(v_c)\bigr],
\]
where occurrences, including repeated occurrences, are sampled uniformly.

\begin{theorem}[Perfect-completeness Label Cover]\label{thm:label-cover}
For every fixed rational $\sigma\in(0,1)$ there are nonempty constant
alphabets and a deterministic polynomial-time reduction from $3$SAT to
explicit Label Cover instances with nonempty occurrence lists such that
satisfiable formulas have value $1$ and unsatisfiable formulas have value
at most $\sigma$.
\end{theorem}

This is the standard consequence of the PCP Theorem and parallel
repetition \cite{AroraEtAl1998,Raz1998}. A quantitative statement is
\cite[Theorem~6.2, full version]{DinurSteurer2014}:
for a fixed number $k$ of repetitions, the output size is $n^{O(k)}$,
the alphabet size is at most $a^k$, and the soundness is at most $\beta_0^k$,
for absolute constants $a>1$ and $0<\beta_0<1$.
Choose a fixed $k$ for which $\beta_0^k\le\sigma$.
The unweighted total-projection formulation is explicitly recorded in
\cite[Definitions~1.1--1.2 and Theorem~1.3]{DKKMS2025}.
The latter source even supplies $d$-to-one maps, a property we do not use.
Here the verifier's random choices specify an explicit list of constraints;
they are not random choices of the reduction.
Unused questions can be removed.

\subsection{Juntas for the sup metric}

A function is a \emph{junta on $S$} if it depends only on the input
coordinates in $S$. The approximating functions below need not be continuous.
We use Austin's continuous junta theorem \cite{Austin2016}, which
extends the dimension-independent coordinate approximation phenomenon
of Friedgut's Boolean junta theorem \cite{Friedgut1998}.

\begin{theorem}[Dimension-independent junta approximation]
\label{thm:junta}
For every $0\le L<\infty$ and $\tau>0$ there is an integer $J(L,\tau)\ge1$
such that every $L$-Lipschitz function
$f:[0,1]^N\to[-1,1]$, for every positive integer $N$, admits a junta
on at most $J(L,\tau)$ coordinates with $L^2$ error less than $\tau$.
The junta may be taken to be a coordinate conditional expectation of $f$.
\end{theorem}

To obtain this formulation from
\cite[Theorem~1.1, arXiv version~4]{Austin2016},
use the usual interval with uniform measure and modulus of continuity
$\omega(t)=Lt$. The interval is compact, connected, and locally connected,
as required there. Austin's theorem gives
$\|f-\E[f\mid X_S]\|_1<\tau^2/2$ with $|S|$ bounded independently of $N$.
Both functions take values in $[-1,1]$, so
\[
 \|f-\E[f\mid X_S]\|_2^2
 \le 2\|f-\E[f\mid X_S]\|_1<\tau^2.
\]
The bound applies to pullbacks of torus functions, since the quotient
$[0,1]^N\to\T^N$ is nonexpanding and preserves the product uniform law.

For a square-integrable function $H$ of independent circle coordinates
$z_1,\ldots,z_s$, let $E_i$ average coordinate $z_i$ alone, and set
\[
 Q_i=\id-E_i,\qquad D_i(H)=\|Q_iH\|_2.
\]
Both $E_i$ and $Q_i$ are orthogonal projections. In particular,
$\|Q_if\|_2\le\|f\|_2$; the constant here is one.

\begin{lemma}[Two dimension-independent bounds]\label{lem:influences}
Fix $0\le L<\infty$ and $\eps>0$. There are $\beta>0$ and an integer $K\ge1$,
independent of $s$, such that every $L$-Lipschitz
$H:\T^s\to[-1,1]$ satisfies:
\begin{enumerate}[label=\textup{(\roman*)}]
\item if $\E H=0$ and $\E H^2>\eps$, then $D_i(H)\ge\beta$ for some $i$;
\item at most $K$ indices satisfy $D_i(H)\ge\beta/2$.
\end{enumerate}
\end{lemma}

\begin{proof}
Let $J=J(L,\sqrt{\eps}/2)$ from Theorem~\ref{thm:junta}.
Choose a set $S$ of at most $J$ coordinates such that
$G=\E[H\mid z_S]$ satisfies $\|H-G\|_2<\sqrt{\eps}/2$.
If $H$ has mean zero and squared norm exceeding $\eps$, then
$G$ has mean zero and, by orthogonality,
$\|G\|_2^2>3\eps/4$, in particular $\|G\|_2^2>\eps/4$.

For completeness, decompose
\[
 H=\sum_{A\subseteq[s]} H_A,\qquad
 H_A=\Bigl(\prod_{i\in A}Q_i\Bigr)
     \Bigl(\prod_{i\notin A}E_i\Bigr)H .
\]
The factors commute and the summands are orthogonal.
Thus $G=\sum_{A\subseteq S}H_A$ and, in the mean-zero case,
\[
 \|G\|_2^2
 =\sum_{\varnothing\ne A\subseteq S}\|H_A\|_2^2
 \le \sum_{i\in S}\sum_{A\ni i}\|H_A\|_2^2
 =\sum_{i\in S}D_i(H)^2.
\]
Consequently $\beta=\sqrt{\eps/(4J)}$ works in (i).

For (ii), use Theorem~\ref{thm:junta} again with error less than
$\beta/2$, and write $G'$ for the resulting junta, on at most
$K=J(L,\beta/2)$ coordinates. If $i$ is not one of these coordinates,
then $Q_iG'=0$, so
\[
 D_i(H)=\|Q_i(H-G')\|_2\le\|H-G'\|_2<\beta/2.
\]
This proves the cardinality bound.
\end{proof}

The first part locates a coordinate when the function has substantial
variance; the second limits how many coordinates can have a smaller
but still fixed amount of dependence. Their simultaneous dimension
independence is what will permit the number of test blocks to be chosen
after both bounds are fixed.
